\scanpage{33}

\deleted{Graphe
\diagram{\begin{tikzcd}[ampersand replacement=\&,column sep=small,row sep=small]
(S,\arrow[d]\&A,\arrow[d]\&R,\arrow[d]\&\sigma)\arrow[d]\&R\subset S\times A\\
\text{sommets}\&\substack{\text{arêtes}\\\text{\unsure{orientées|orientée}}}\&\substack{\text{relation}\\\text{\unsure{incidence}}}\&\text{symétrie}\&
\end{tikzcd}}}
\deleted{Graphe $S$ : ensemble $S$ avec une partie $A$ de $S\times S$, stable par symétrie}

Graphe
\[
 (S,A,A'\xrightarrow{\varphi=(\varphi_1,\varphi_2)=(\varphi_1,\varphi_1\circ\sigma)}S\times S,\sigma)
 \Longleftrightarrow (S,R\subset A\times A,L\longrightarrow A,\sigma)
\]
\diagram{\begin{tikzcd}[ampersand replacement=\&,row sep=small,column sep=small]
 S\arrow[d]\&A\arrow[d]\&A'\arrow[d,bend right=15]\&\varphi\arrow[d]\&\sigma\arrow[d]\\
 \text{sommets}\&\substack{\text{arêtes}\\\text{orientées}}\&\substack{[\text{arêtes}\\\text{non circulaires isolées}]}\&\substack{\deleted{\text{\unsure{symétrie}}}\\\deleted{\text{\unsure{départ}}}}\&\substack{\text{symétrie}\\\text{dans }A}
\end{tikzcd}}
\inserted{$R$ stable par sym.\qquad $L$ arêtes \unsure{boucles|bouclées}}

$G$ groupe des symétries du graphe \qquad $G\simeq\mathrm{Perm}^{\unsure{\prime}}(S)\times\mathrm{Perm}(A)$

Supposons que $G$ est transitif sur $A$ et sur $S$.

Alors a) \deleted{tous les \unsure{pts}} ou bien tous les \deleted{points}\inserted{sommets} sont isolés (et les arêtes sont circulaires isolées) donc le graphe est juste formé de deux ensembles $S,A$, $A'=\varnothing$, et un autom d'ordre $2$ dans $A$ \deleted{\unsure{commutant à}}\inserted{sans pt fixe}. \deleted{l'action de $G$. Donc}
\[
 \mathrm{Aut}\,\Gamma=\mathrm{Perm}(S)\times\mathrm{Aut}(A,\sigma)
\]
% slanted
\margin{tous les sommets isolés}

\margin{Si $A_{1}=A/\sigma$, on trouve que $\mathrm{Aut}(A,\sigma)\simeq\mathrm{Perm}(A_1).(\mathbb Z/2\mathbb Z)^{A_1}$ Ceci doit \unsure{servir|revenir|semi} d'\unsure{int|intro|direct}}

\deleted{[b) ou bien \inserted{il y a des sommets non isolés, donc} tous les sommets sont non isolés, \deleted{et il y a des points \unsure{et donc}} Dans ce cas, il y a des arêtes non \inserted{circulaires} isolées, \deleted{ou bien} toutes les arêtes sont non \inserted{circulaires} isolées [car il y en a non isolées]

% slanted
\margin{il existe des arêtes non isolées circulaires}

\deleted{Il est possible} Ou bien toutes les arêtes sont circulaires, donc $\varphi$ se réduit à une application \uline{surjective}
\[
 \varphi:A\longrightarrow S
\]
compatible avec $\sigma$ $[\varphi\sigma x=\varphi x]$

\deleted{On aura $A$} Choisissons $a\in A$, donc $A\simeq G/G_a^{\inserted{=H}}$ (groupe d'isotropie) et si $s=\varphi(a)$, $S\simeq G/G_s^{\inserted{=H'}}$, $\underset{\unsure{H'}}{G_s}\supset\unsure{A|H|G_a}$ : L'application $\sigma:A\longrightarrow A$ correspond à la donnée d'un élément de $\dot\sigma\in\mathrm{Norm}(G_a)/G_a$, $\dot\sigma$ d'ordre $2$ ; le fait que $\unsure{\rho|\varphi}$ \deleted{\unsure{commute}}\inserted{est compatible} avec $\sigma$ signifie que $\sigma\in G_s=H'$ ; le fait que $\sigma$ opère sans pts fixes signifie que $\sigma\notin H\deleted{\unsure{\bigcup_{g\in G}gHg^{-1}}}$ ; donc on a
\[
 \left\{\begin{aligned}
 &\boxed{G\supset H'\supset H}\quad\text{et}\\
 &\deleted{\sigma\in H'\cap\mathrm{Norm}_G(H)\cap\complement(\unsure{\bigcup_{g\in G}gHg^{-1}})H}\\
 &\boxed{\dot\sigma\in(H'\cap\mathrm{Norm}_G(H))/H\qquad\dot\sigma^2=1,\ \dot\sigma\ne1}
 \end{aligned}\right.
\]
Le fait que $G$ opère fidèlement \inserted{sur $\Gamma$} signifie que
\[
 \bigcap_{g\in G}gHg^{-1}=\{1\}
\]
Si on a une telle situation \hfill ]}
% archive sheet number: 32
